% ============================================================================
% 01_intro.tex — auction-v3 (screening reframe)
% THIRD PASS per Kehang (2026-07-19), styled on skills/writeup/example.pdf
% (Manning-Horton "General Social Agents" intro): SHORT paragraphs, one job
% each; the main-result preview is its OWN paragraph (P5), phrased
% mechanism-generically so it does not lean on the auction scenario, which
% is introduced only afterwards (P6, opening with "This paper needs one
% such mechanism..."). Order:
%   P1 phenomenon+stakes | P2 design response + search bottleneck |
%   P3 opportunity (LLM agents) | P4 what we do: ordinal screening + RQ |
%   P5 MAIN-RESULT PREVIEW (own paragraph, general) | P6 scenario (SPSB) |
%   P7 clock contrast | P8 theory | P9 design space | P10 discipline+cost |
%   P11 contributions+scope | P12 benchmark layer | P13 validation numbers |
%   P14 opposite errors + cardinal boundary | P15 map results (robust) |
%   P16 flags+predictions+portability | P17 traces | P18 roadmap.
% Related literature woven inline where the content appears (no standalone
% review). Canon numbers per contents_v2/CONVENTIONS.md §3.
% ============================================================================
\section{Introduction}\label{sec:intro}

% --- P1: The phenomenon and its stakes --------------------------------------
In many market mechanisms, participants do not play the strategy that theory identifies as optimal---even when that strategy is simple, safe, and known to the designer. Auction bidders pay more for items than the items are worth to them \citep{kagel1995individual}. Applicants in deployed school-choice and labor-market clearinghouses submit dominated preference rankings even where reporting honestly is provably optimal \citep{rees2018suboptimal, roth1982economics}. Decades of laboratory evidence show these deviations to be systematic rather than noise \citep{kagel2020handbook}---less a matter of exotic preferences than of the computational difficulty of recognizing what the mechanism makes safe \citep{oprea2024decisions}. The stakes are the mechanism's own objectives: dominated play erodes the efficiency, revenue, and fairness properties the rules were chosen to deliver.

% --- P2: The design response, and the search bottleneck ---------------------
The design response is to make correct play easy, and a well-developed theory now characterizes which mechanisms are cognitively simple \citep{li2017obviously, li2024designing}. But for a practitioner the question is unavoidably empirical: among the many concrete choices of implementation, wording, and decision support, which ones preserve correct play in \emph{their} market? The standard instrument for that question---the incentivized human experiment---is slow and expensive enough that only a handful of design variants are ever tested, each chosen well in advance. Simplicity has therefore remained a subject for theory and isolated experiments rather than a routine engineering practice.

% --- P3: The new opportunity: LLM agents ------------------------------------
The new opportunity is that a different kind of experimental participant has arrived. Large language models (LLMs) reproduce qualitative patterns from classic behavioral experiments \citep{horton2023large, aher2023using} and support automated experimentation at scale \citep{park2023generative, manning2024automated}. They have also begun to transact in real markets as bidding and negotiating agents in their own right---studied for their collusive tendencies \citep{fish2024algorithmic}, their feasibility as bidders \citep{tsuchihashi2023homo, lu2025claude}, and their trainability as auction participants \citep{chen2023put}---so their behavior under a mechanism's rules is becoming designer-relevant in itself. A designer can now run a mechanism experiment on artificial participants in minutes, for dollars, with reproducible seeds.

% --- P4: What we do: ordinal screening --------------------------------------
This paper develops and validates a way to use that opportunity that does not require the agents to be human-like: \emph{ordinal screening}. The obvious alternative---treat the simulated population as a stand-in for human subjects and read levels off it---asks the instrument for more than it can give: LLM agents do not reproduce human levels, they do not err in the human direction, and although per-task calibration can narrow distributional gaps \citep{corrigan2025usefulness}, we show that no tuning of the agents is portable across mechanisms (Appendix~\ref{app:cardinal}). Under ordinal screening, the designer asks the panel only which design is behaviorally \emph{easier}---directions and orderings, never levels---with the design space fixed ex ante by theory, a cross-model robustness rule deciding what counts as a finding, and scarce human experiments reserved for the cells the screen elevates. The question the paper answers is: \emph{can LLM agents recover robust ordinal patterns in human behavior across mechanisms, and can those ordinal signals be used to screen theory-guided mechanism designs?}

% --- P5: Main-result preview (own paragraph; mechanism-generic) -------------
Both answers are yes. On every mechanism comparison for which published human evidence exists---which of two formats is behaviorally harder, whether a simpler implementation of the same rule improves play, whether restating the rules moves behavior---a panel of four LLM families recovers the direction of the human finding, family by family (Figure~\ref{fig:ordinal}). The agreement is not inherited from mimicry: the two populations' characteristic errors point in \emph{opposite} directions, so a panel that merely reproduced human data from training would get the comparisons wrong, and no tuning we searched makes the agents match human behavior in levels. Agreement is ordinal, not cardinal---and ordinal is the signal design screening needs. Used this way, the panel populates a theory-organized design map at roughly \$10 per cell, turning simplicity from a property that theory characterizes into a quantity a practitioner can engineer.

% --- P6: The concrete scenario: the SPSB auction -----------------------------
This paper needs one such mechanism as its concrete scenario, and we use the canonical one throughout: the second-price sealed-bid (SPSB) auction. Each bidder privately submits one bid; the highest bidder wins but pays the \emph{second}-highest bid. Bidding exactly what the item is worth to you is optimal no matter what the other bidders do---your bid determines only \emph{whether} you win, never \emph{what} you pay \citep{vickrey1961counterspeculation}. Economists call such a strategy \emph{dominant}, and mechanisms that make truthful reporting dominant \emph{strategy-proof}. Human bidders violate the prescription anyway, most commonly by bidding \emph{above} value, persistently and expensively \citep{kagel1993independent, kagel1995individual}.

% --- P7: The clock contrast, and why auctions --------------------------------
The failure is implementation-dependent, which is what makes it a design problem rather than a fixed cost of human nature. Run the same allocation rule as an ascending clock---the price rises step by step, each bidder chooses \emph{stay} or \emph{exit} at each price, and the last bidder standing wins at the runner-up's exit price---and human subjects find the dominant strategy far more often \citep{li2017obviously, breitmoser2022obviousness}. Same economics, different interface, different behavior. We choose auctions as the scenario because they pair the sharpest dominant-strategy predictions with the deepest human experimental record to validate against; nothing in the method is specific to them.

% --- P8: The theory behind simplicity -----------------------------------------
A body of theory explains why implementation matters, and it supplies the coordinates of our design space. \citet{li2017obviously} formalizes \emph{obvious strategy-proofness} (OSP): in the clock, but not in the sealed bid, verifying that truth-telling is best requires no hypothetical, case-by-case reasoning about what opponents might have bid. Related work characterizes when optimal play requires no beliefs about others at all \citep{borgers2019strategically}, when it requires no multi-step foresight \citep{pycia2023theory}, and how a mechanism's cognitive burden decomposes into contingent reasoning, forward planning, and belief formation \citep{li2024designing}. None of this is auction-specific: the same notions classify matching and assignment mechanisms \citep{ashlagi2018stable}, and we demonstrate the instrument's portability by replicating the paper's central description finding under student-proposing deferred acceptance (Appendix~\ref{app:da}). What the theory does not do is rank the concrete choices a designer faces; it constrains the design space and signs its cells ex ante, and the rest is measurement.

% --- P9: The design space ------------------------------------------------------
Section~\ref{sec:theory} turns that theory into a design map with three dimensions, each holding the allocation rule fixed. \emph{Implementation}: replace the one-shot sealed bid with the OSP ascending clock \citep{li2017obviously, pycia2023theory}. \emph{Description}: keep the sealed bid and change only the words---a menu-style restatement of the outcome function, which improves measured understanding while leaving human bidding essentially unchanged \citep{gonczarowski2023strategyproofness, gonczarowski2024describing}; a one-sentence \emph{safety} description stating the invariant that makes truthful play safe (``your bid determines \emph{if} you win, not \emph{what} you pay''), never tested in a human auction, though property-exposing descriptions help and mechanics-only descriptions backfire in matching mechanisms \citep{katuscak2024simpler, guillen2018effectiveness}; and a framing of the sealed bid as an automatic exit threshold on a hypothetical clock, which raises human truthful bidding \citep{breitmoser2022obviousness}. \emph{Reasoning support}: hold the mechanism and its description fixed and scaffold the participant's reasoning---walk it through the payoff consequences of each outcome, or prompt it to think about opponents, an axis strategy-proofness makes irrelevant and therefore a built-in diagnostic \citep{li2024designing, borgers2019strategically}. Direct advice to play truthfully is excluded by construction: it works on humans, but it reveals the answer rather than exposing why the answer is safe \citep{masuda2022net}.

% --- P10: Discipline and the cost of search ------------------------------------
Two features discipline the exercise. First, every cell's predicted sign is fixed ex ante from the theory, and roughly half the cells are predicted to be inert or to backfire---a pattern opportunistic prompt search would neither produce nor publish. Second, the economics of search are on the table: the laboratory test of obvious strategy-proofness engaged 548 subjects with documented payments around \$5{,}400 for its main experiment alone \citep{li2017obviously}, and the leading description experiments recruited 542 subjects at \$15--26 each \citep{gonczarowski2024describing}, while the panel sweeps the full design grid for under \$400.

% --- P11: Contributions and scope ----------------------------------------------
The paper contributes a validated screening instrument and the first populated map of this design space. The validation contribution establishes, comparison by comparison, that the instrument's ordinal readings agree with the human experimental record wherever that record exists, under a pre-committed rule: a result enters the main text only if its direction is consistent across all four prespecified model families. The discovery contribution fills the map, classifies levers into substitutes, complements, and anti-levers, and converts every cell that no human experiment has visited into a concrete, falsifiable prediction with a costing. Throughout, the panel is evaluated as an ordinal design-screening tool, not as a calibrated model of the average human bidder. Our question is distinct from designing the internal auctions of LLM-powered platforms \citep{duetting2024mechanism} and from automated mechanism design, which searches over allocation rules themselves \citep{dutting2024optimal, wang2024gemnet}: we hold the rule fixed and engineer the interface a bounded reasoner meets.

% --- P12: The human benchmark layer --------------------------------------------
Screening requires a benchmark layer, and ours is built from the human experimental record. Because micro-level bid data from the classic experiments are rarely available, we reconstruct human bidding distributions by moment matching: a transparent mixture model calibrated, source by source, to the aggregate statistics each paper reports \citep{kagel1993independent, li2017obviously, breitmoser2022obviousness, gonczarowski2023strategyproofness}. A single deviation metric---the scaled mean absolute deviation (SMAD) of actions from the theoretical benchmark---places humans and LLM agents on the same scale, and a \emph{preservation index} expresses every design lever's effect as the fraction of a population's own baseline deviation that the lever removes. The reconstructions are held to a strict epistemic standard, stated once and enforced throughout: comparisons to reconstructed human behavior are made at the level of rankings, directions, and comparative statics only; levels and significance tests are reserved for theory benchmarks.

% --- P13: Validation results, with numbers --------------------------------------
The validation results are collected in one exhibit (Figure~\ref{fig:ordinal}), and every row shows its per-family direction. The panel---GPT-4o, Claude 3.5 Haiku, Gemini 2.0 Flash, Gemma 27B---reproduces the human difficulty ordering of payment rules: first-price bidding, where equilibrium requires strategic shading, is harder than second-price bidding, where truth-telling is dominant, in all four families (GPT-4o: SMAD 27.6\% versus 8.7\%, against a human 24.8\% versus 5.7\%). It reproduces the implementation effect: the ascending clock moves every family toward truthful play ($p<0.001$ in each), as it moves humans. And it reproduces the description null: menu-style restatements, which leave human bidding essentially unchanged, have no consistent effect on the panel either.

% --- P14: Opposite errors, and the instrument's boundary -------------------------
The agreement is more informative than it first appears, because the two populations make \emph{opposite} baseline mistakes. Humans overbid in second-price auctions---67.2\% of reconstructed bids exceed value---while the models underbid: 74.1\% of GPT-4o's baseline bids fall below value. That the same design levers repair opposite errors in both populations indicates that the levers act on the strategic problem itself---the recognizability of the dominant strategy---rather than on any psychology particular to humans. It also rules out the deflationary reading that the panel parrots the human experimental literature from its training data: a parroting population would overbid. The same divergence marks the instrument's boundary, measured in Appendix~\ref{app:cardinal}: no tuning of the agents---temperature, personas, framings, explanations---achieves level-and-direction agreement across mechanisms, and the one knob that restores the human direction of error in second-price auctions induces dominated bids in first price, where humans almost never overbid.

% --- P15: Design-map results: the four robust findings ---------------------------
The discovery results fill the map (Section~\ref{sec:map}, Table~\ref{tab:design-map}), and four findings pass the cross-family robustness rule. The OSP clock implementation is the strongest lever, compressing deviations toward zero in every family and eliminating the catastrophic tail of severe underbids. A one-sentence safety description---honest text stating \emph{why} truthful bidding is safe---recovers roughly half of the clock's benefit in every family, at the cost of a sentence. A payoff-tree scaffold, which lays out what happens in each contingency without recommending an action, delivers a comparable improvement. And belief-formation scaffolds backfire, pushing bids further from value, exactly as the theory's diagnostic logic predicts for agents that have not internalized strategy-proofness.

% --- P16: Flags, predictions, portability -----------------------------------------
The rest of the map is flagged or forward-looking. Two cells are honestly heterogeneous and reported as such: the menu restatement is null on average but sign-mixed underneath, and a clock-\emph{framing} description helps three families while harming the fourth---a paraphrase sensitivity we quantify as a warning against single-wording inferences (Appendix~\ref{app:heterogeneity}). Cells that no human experiment has visited become predictions with directions and costings---most cheaply, that the safety sentence will shift human bidding toward value where menu restatements have not, testable in a single incentivized session. And the sharpest finding travels: in the school-choice matching domain, a description carrying the safety invariant nearly eliminates misreporting while a description restating only the mechanics makes it strictly worse (Appendix~\ref{app:da}).

% --- P17: The traces -----------------------------------------------------------------
The instrument also records \emph{why} agents say they act as they do, which yields a methodological result for anyone auditing an interface change by asking participants to explain themselves. Every bid arrives with a stated plan---a short written explanation of the intended strategy, such as ``I will bid slightly below my value to keep a profit margin.'' Coding all 21{,}990 plans with a keyword dictionary aligned to the theory's cognitive categories, we find that stated reasoning and realized bids never move together across our levers (Section~\ref{sec:traces}): the safety description and the payoff tree move bids while leaving stated plans unchanged; a worst-case scaffold shifts stated plans while leaving bids unchanged; the menu restatement moves neither. Effective designs changed what agents \emph{did}, not what they \emph{said}---so screening, synthetic or human, should be evaluated on realized play rather than elicited understanding.

% --- P18: Roadmap ---------------------------------------------------------------------
The paper proceeds as follows. Section~\ref{sec:theory} develops the simplicity framework---first for a general mechanism, then in the auction scenario---and presents the design map with its ex-ante predictions. Section~\ref{sec:validation} describes the testbed and reports the ordinal validation. Section~\ref{sec:map} populates the design map, reports the robust lever results and the flagged heterogeneous ones, and registers the interaction cells for future runs. Section~\ref{sec:traces} analyzes stated plans against realized bids. Section~\ref{sec:discussion} collects design implications, the screening workflow and its cost, robustness and scope (including frontier reasoning models, which solve our baseline formats yet remain vulnerable to presentation---Appendix~\ref{app:ablations}), limitations, and the predictions for human experiments, and concludes. Appendices collect procedures, prompts, the human reconstructions, the cardinal-calibration analysis, model-specific heterogeneity, and the school-choice matching domain in which the description finding replicates (Appendix~\ref{app:da}).
